% =====================================================================
%  2bat-ud1-8.tex  —  SESSIÓ 8: Quantes hores de cada professional?
%  (sense preàmbul: s'inclou des de main.tex)
%  Criteris: C3 (producte) i C4 (interpretar). Estudi de cas (Activitat 3).
% =====================================================================
\horatitol{Sessió 8 — Quantes hores de cada professional?}%
{Un cas real: multiplicar dues taules per saber les hores de dedicació que genera cada usuari d'un servei.}

\subsection{Idea clau}

A la sessió anterior vam aprendre \emph{com} es fa el producte de matrices (fila
per columna). Ara el posem a treballar en un cas real de planificació. Imagina un
servei d'atenció a les persones on cada usuari necessita un nombre determinat de
sessions de diferents tipus (educativa, d'acompanyament, administrativa), i cada
tipus de sessió té una durada fixa en hores.

Si volem saber \textbf{quantes hores de dedicació} genera cada usuari, multipliquem
la matriu de sessions per la matriu d'hores. Cada número del resultat ens diu la
\emph{càrrega de treball total} que suposa un usuari concret per al servei.

\begin{idea}
\textbf{Llegir el resultat d'un producte}
\begin{itemize}[leftmargin=1.4em, itemsep=2pt]
  \item Si una fila són les \textbf{sessions} d'un usuari i una columna són les
        \textbf{hores} per sessió, el producte fila $\times$ columna és el
        \textbf{total d'hores} d'aquell usuari.
  \item Cada element del resultat té un significat concret: no és un número
        abstracte, és \emph{dedicació real} que cal preveure.
  \item Aquesta eina serveix per \textbf{dimensionar la plantilla}: quins usuaris
        generen més càrrega i quantes hores de professional caldrà en total.
\end{itemize}
\end{idea}

\subsection{Exemples}

\begin{exemple}[les hores totals d'un usuari]
Un usuari necessita, cada setmana, $2$ sessions educatives, $1$ d'acompanyament i
$3$ administratives. Les durades respectives són $2$, $3$ i $1$ hores. Les hores
totals de dedicació són el producte fila per columna:
\[
\begin{pmatrix} 2 & 1 & 3 \end{pmatrix}\cdot\begin{pmatrix} 2 \\ 3 \\ 1 \end{pmatrix}
= 2\per 2 + 1\per 3 + 3\per 1 = 4 + 3 + 3 = 10 \text{ hores}.
\]
Aquest usuari genera $10$ hores de dedicació setmanal per al servei.
\end{exemple}

\begin{exemple}[dos usuaris alhora]
Si afegim un segon usuari (que necessita $1$ educativa, $2$ d'acompanyament i $1$
administrativa), posem cada usuari en una fila:
\[
S=\begin{pmatrix} 2 & 1 & 3 \\ 1 & 2 & 1 \end{pmatrix},\qquad
H=\begin{pmatrix} 2 \\ 3 \\ 1 \end{pmatrix}.
\]
El producte dona les hores totals de cada usuari:
\[
S\cdot H=\begin{pmatrix} 2\per 2+1\per 3+3\per 1 \\ 1\per 2+2\per 3+1\per 1 \end{pmatrix}
        =\begin{pmatrix} 10 \\ 9 \end{pmatrix}.
\]
El primer usuari genera $10$ hores i el segon, $9$.
\end{exemple}

\subsection{Exercicis resolts}

\begin{exresolt}[planificar la dedicació de tres usuaris]
Un servei atén tres usuaris. La matriu recull quantes sessions de cada tipus
(columnes: educativa, acompanyament, administrativa) necessita cada usuari (files):
\[
S=\begin{pmatrix} 3 & 2 & 1 \\ 1 & 4 & 2 \\ 2 & 1 & 3 \end{pmatrix}.
\]
Les durades de cada tipus de sessió són $2$, $1$ i $1$ hores:
\[
H=\begin{pmatrix} 2 \\ 1 \\ 1 \end{pmatrix}.
\]
\begin{enumerate}[label=\alph*), leftmargin=1.6em, topsep=2pt]
  \item Calcula les hores totals de dedicació de cada usuari, $S\cdot H$.
  \item Quin usuari genera més càrrega de treball?
  \item Quantes hores de professional caldrà preveure en total per als tres usuaris?
\end{enumerate}
\solucio
a) Multipliquem fila per columna:
\[
S\cdot H=\begin{pmatrix} 3\per 2+2\per 1+1\per 1 \\ 1\per 2+4\per 1+2\per 1 \\ 2\per 2+1\per 1+3\per 1 \end{pmatrix}
        =\begin{pmatrix} 9 \\ 8 \\ 8 \end{pmatrix}.
\]
b) El primer usuari genera més càrrega: $9$ hores (els altres dos, $8$ cadascun).

c) En total cal preveure $9 + 8 + 8 = 25$ hores de professional a la setmana.
\resposta{$S\cdot H=\begin{pmatrix} 9 \\ 8 \\ 8 \end{pmatrix}$; \;el primer usuari genera més càrrega ($9$ h); \;en total, $25$ hores setmanals.}
\end{exresolt}

\begin{exresolt}[hores i cost alhora]
La matriu de sessions de dos usuaris (files) en dos tipus de sessió (columnes) és
\[
S=\begin{pmatrix} 3 & 2 \\ 1 & 4 \end{pmatrix}.
\]
Cada tipus de sessió té una durada (hores) i un cost (euros):
\[
D=\begin{pmatrix} 2 & 10 \\ 1 & 8 \end{pmatrix}.
\]
Calcula $S\cdot D$ i interpreta cada columna del resultat.
\solucio
Multipliquem fila per columna:
\[
S\cdot D=\begin{pmatrix} 3\per 2+2\per 1 & 3\per 10+2\per 8 \\ 1\per 2+4\per 1 & 1\per 10+4\per 8 \end{pmatrix}
        =\begin{pmatrix} 8 & 46 \\ 6 & 42 \end{pmatrix}.
\]
La primera columna són les \textbf{hores} totals de cada usuari ($8$ i $6$); la
segona columna és el \textbf{cost} total en euros ($46\eur$ i $42\eur$).
\resposta{$S\cdot D=\begin{pmatrix} 8 & 46 \\ 6 & 42 \end{pmatrix}$: \;usuari 1, $8$ h i $46\eur$; \;usuari 2, $6$ h i $42\eur$.}
\end{exresolt}

\subsection{Exercicis per fer}

\begin{exercicis}
\begin{exlist}
  \item Un usuari necessita $4$ sessions educatives, $2$ d'acompanyament i $1$
        administrativa, amb durades de $2$, $1$ i $3$ hores respectivament. Escriu
        el producte fila per columna i calcula les hores totals de dedicació.
  \item Un servei atén dos usuaris amb tres tipus de sessió:
        \[ S=\begin{pmatrix} 2 & 3 & 1 \\ 4 & 1 & 2 \end{pmatrix},\qquad
           H=\begin{pmatrix} 1 \\ 2 \\ 2 \end{pmatrix}. \]
        Calcula $S\cdot H$ i digues quin usuari genera més hores de dedicació.
  \item Tres usuaris necessiten dos tipus de sessió (educativa i d'acompanyament):
        \[ S=\begin{pmatrix} 3 & 1 \\ 2 & 2 \\ 1 & 4 \end{pmatrix},\qquad
           H=\begin{pmatrix} 2 \\ 1 \end{pmatrix}. \]
        Troba les hores de cada usuari i el total d'hores que cal preveure.
  \item Amb la matriu de sessions $S=\begin{pmatrix} 2 & 1 \\ 3 & 2 \end{pmatrix}$
        i la taula de durada i cost $D=\begin{pmatrix} 2 & 12 \\ 1 & 9 \end{pmatrix}$,
        calcula $S\cdot D$ i digues quantes hores i quin cost genera cada usuari.
  \item Un casal té quatre usuaris i dos tipus de taller (manualitats i esport):
        \[ S=\begin{pmatrix} 2 & 1 \\ 1 & 3 \\ 0 & 2 \\ 3 & 1 \end{pmatrix},\qquad
           H=\begin{pmatrix} 1 \\ 2 \end{pmatrix}. \]
        Calcula les hores de cada usuari i digues quin en genera menys.
  \item \textbf{(Raonament)} En el cas de l'exercici 3, un company calcula
        $H\cdot S$ en lloc de $S\cdot H$ i diu que «és igual». Sense fer tot el
        càlcul, explica per què l'ordre importa i si $H\cdot S$ es pot fer (mira
        les dimensions: $S$ és $3\times 2$ i $H$ és $2\times 1$).
\end{exlist}
\end{exercicis}

% ---------------------------------------------------------------------
\fullsolucions{Sessió 8}
\begin{solucions}
\begin{sollist}
  \item $\begin{pmatrix} 4 & 2 & 1 \end{pmatrix}\cdot\begin{pmatrix} 2 \\ 1 \\ 3 \end{pmatrix}
        =4\per 2+2\per 1+1\per 3=8+2+3=13$ hores.
  \item $S\cdot H=\begin{pmatrix} 10 \\ 10 \end{pmatrix}$;
        \;els dos usuaris generen les mateixes hores ($10$ cadascun).
  \item $S\cdot H=\begin{pmatrix} 7 \\ 6 \\ 6 \end{pmatrix}$;
        \;cal preveure $7+6+6=19$ hores en total.
  \item $S\cdot D=\begin{pmatrix} 5 & 33 \\ 8 & 54 \end{pmatrix}$;
        \;usuari 1: $5$ h i $33\eur$; usuari 2: $8$ h i $54\eur$.
  \item $S\cdot H=\begin{pmatrix} 4 \\ 7 \\ 4 \\ 5 \end{pmatrix}$;
        \;el primer i el tercer usuari en generen menys ($4$ hores cadascun).
  \item L'ordre importa perquè en el producte de matrices, en general,
        $S\cdot H\neq H\cdot S$. A més, $H\cdot S$ \emph{no} es pot fer: $H$ és
        $2\times 1$ i $S$ és $3\times 2$, i les columnes de $H$ ($1$) no coincideixen
        amb les files de $S$ ($3$). Per tant $H\cdot S$ no és compatible.
\end{sollist}
\end{solucions}
